% !PRINTED 740
\weberpage{740}

%<<CONT
kommt, es sei $a_{1,1}$ diese ungerade Zahl, dann bestimmen wir
eine zweite Lösung:
\[
0, a_{2,2}, a_{3,2}, \ldots, a_{\nu-1,2},
\]
in der wieder eine ungerade Zahl, etwa $a_{2,2}$, vorkommt. Dann
bestimmen wir eine dritte Lösung:
\[
0, 0, a_{3,3}, \ldots, a_{\nu-1,3},
\]
und fahren mit Nullsetzen so lange fort, als die Anzahl der
übrig bleibenden Unbekannten noch grösser ist, als die Anzahl
der Gleichungen.

Im letzten Systeme werden also $\nu'-1$ Nullen vorkommen,
damit $\nu'$ Unbekannte übrigbleiben.

Diese Zahlenreihen $a_{\nu,k}$ können wir für das System (20)
setzen, denn darin reducirt sich die Determinante
\weq{21}{\Sigma \pm a_{1,1}\, a_{2,2} \ldots a_{\nu',\nu'}}
auf das Diagonalglied $a_{1,1}\, a_{2,2} \ldots a_{\nu',\nu'}$, dessen Factoren alle ungerade
 sind.

Ist nun die Forderung des Satzes 2. erfüllt, so lässt sich
die Matrix (20) durch Hinzufügung von $\nu'-1$ Zeilen zu einer
Determinante von $\nu-1$ Reihen ergänzen, die gleichfalls einen
ungeraden ganzzahligen Werth hat:
\weq{22}{a = \Sigma \pm a_{1,1}\, a_{2,2} \ldots a_{\nu',\nu'}\, a_{\nu'+1,\nu'+1} \ldots a_{\nu-1,\nu-1}.}

Man braucht nur, wenn die Determinante (21) als ungerade
vorausgesetzt wird, in den hinzugefügten Zeilen für die Elemente
lauter Nullen zu setzen, ausgenommen die in der Diagonalreihe
 stehenden, $a_{\nu'+1,\nu'+1} \ldots a_{\nu-1,\nu-1}$, die gleich 1 genommen
werden können. Dann erhält $a$ denselben Werth wie die Determinante
 (21).

Wenn wir aber jetzt nach dem Multiplicationssatze der
Determinanten das Product $a M$ bilden, so ergiebt sich eine
Determinante, in der die $\nu-1$ Summen
\[
\begin{aligned}
& \sum\limits_{1,\, \nu-1}^{i} a_{i,s}\, m_{k,i}, \qquad \sum\limits_{1,\, \nu-1}^{i} a_{i,s}\, M_{k,i} \\
& k = 1, 2, \ldots \nu'-1, \qquad k = 0, 1, 2 \ldots \nu'-1
\end{aligned}
\]
für jeden Werth $s = 1, 2, \ldots, \nu'$ eine Reihe bilden, und in der
also $\nu'$ Reihen vorkommen, deren Elemente alle durch 2 theilbar
 sind.
